\section{Word Ladder II}
\label{sec:graph-word-ladder-2}

\problemstatement{Given $\textit{beginWord}$, $\textit{endWord}$, and
$\textit{wordList}$, return \textbf{every} shortest transformation
sequence from $\textit{beginWord}$ to $\textit{endWord}$ (not just the
length, as in Section~\ref{sec:graph-word-ladder-1}).}

\begin{patternbox}
\emph{``Return every shortest path, not just its length''} needs the
same BFS as Word Ladder I, but with one critical addition: a
\textbf{parent map} recording, for every word, every word at the
\emph{previous level} that could reach it -- the same ``print the
path'' machinery used by reconstruction problems throughout this book
series (e.g.\ the DP on LIS chapter's Print LIS): \emph{measuring} an
optimum and \emph{reconstructing} every optimal path are different
problems, and the second always needs extra bookkeeping kept during
the first.
\end{patternbox}

\subsection*{The same-level subtlety}

A naive BFS removes a word from the ``available'' set the instant it
is reached, to avoid revisiting it. But for \textbf{counting all
shortest paths}, two different words at the \emph{same} level might
both be able to reach the same word at the \emph{next} level -- both
of those connections are valid parts of some shortest path and must be
recorded. So word removal must be \textbf{deferred until an entire
level finishes processing}: only after every word at the current level
has tried every substitution should the newly-discovered next-level
words be removed from the available set.

\subsection*{Algorithm}

\begin{enumerate}
  \item BFS level by level. For each word in the \textbf{current}
        level, generate every single-character substitution; for each
        one found in the available word set, record the current word
        as one of its parents, and collect it into the \textbf{next}
        level's word set (without removing it from the global
        available set yet).
  \item After the whole current level is processed, remove every
        newly-collected next-level word from the global available set
        (the deferred removal), then advance to the next level.
  \item Stop as soon as $\textit{endWord}$ is found among the
        newly-collected words for some level (the shortest distance is
        confirmed, and no longer-path could still be shortest).
  \item Backtrack from $\textit{endWord}$ via the parent map, using
        DFS, collecting every path back to $\textit{beginWord}$, then
        reverse each collected path.
\end{enumerate}

\subsection*{C++ implementation}

\begin{lstlisting}
#include <bits/stdc++.h>
using namespace std;

void backtrack(string word, string& beginWord,
               unordered_map<string, vector<string>>& parents,
               vector<string>& path, vector<vector<string>>& result) {
    path.push_back(word);
    if (word == beginWord) {
        vector<string> reversed(path.rbegin(), path.rend());
        result.push_back(reversed);
    } else {
        for (string& parent : parents[word]) {
            backtrack(parent, beginWord, parents, path, result);
        }
    }
    path.pop_back();
}

vector<vector<string>> findLadders(string beginWord, string endWord,
                                    vector<string>& wordList) {
    unordered_set<string> available(wordList.begin(), wordList.end());
    if (available.find(endWord) == available.end()) return {};

    unordered_map<string, vector<string>> parents;
    unordered_set<string> currentLevel = {beginWord};
    available.erase(beginWord);
    bool found = false;

    while (!currentLevel.empty() && !found) {
        unordered_set<string> nextLevel;

        for (const string& word : currentLevel) {
            string candidate = word;
            for (int i = 0; i < (int)candidate.size(); i++) {
                char original = candidate[i];
                for (char c = 'a'; c <= 'z'; c++) {
                    if (c == original) continue;
                    candidate[i] = c;
                    if (available.find(candidate) != available.end()) {
                        nextLevel.insert(candidate);
                        parents[candidate].push_back(word);
                        if (candidate == endWord) found = true;
                    }
                }
                candidate[i] = original;
            }
        }
        for (const string& word : nextLevel) available.erase(word);
        currentLevel = nextLevel;
    }

    vector<vector<string>> result;
    if (found) {
        vector<string> path;
        backtrack(endWord, beginWord, parents, path, result);
    }
    return result;
}

int main() {
    string beginWord = "hit", endWord = "cog";
    vector<string> wordList = {"hot", "dot", "dog", "lot", "log", "cog"};
    vector<vector<string>> ladders = findLadders(beginWord, endWord, wordList);
    for (auto& path : ladders) {
        for (string& w : path) cout << w << " ";
        cout << endl;
    }
    // hit hot dot dog cog
    // hit hot lot log cog
    return 0;
}
\end{lstlisting}

\subsection*{Dry run}

\dryrun{Same input as Section~\ref{sec:graph-word-ladder-1}'s example.}

Level $0$: $\{\texttt{hit}\}$. Level $1$: $\{\texttt{hot}\}$
(\texttt{hit}'s only valid substitution), $\textit{parents}[\texttt{hot}]=[\texttt{hit}]$.
Level $2$: from \texttt{hot}, both \texttt{dot} and \texttt{lot} match
-- collected together as $\{\texttt{dot},\texttt{lot}\}$,
$\textit{parents}[\texttt{dot}]=[\texttt{hot}]$,
$\textit{parents}[\texttt{lot}]=[\texttt{hot}]$ (neither is removed
from \textit{available} until \emph{both} are found, which is exactly
the deferred-removal rule). Level $3$: from \texttt{dot}, \texttt{dog}
matches; from \texttt{lot}, \texttt{log} matches --
$\{\texttt{dog},\texttt{log}\}$. Level $4$: from \texttt{dog},
\texttt{cog} matches ($\textit{found}=\texttt{true}$); from
\texttt{log}, \texttt{cog} \emph{also} matches (still available, since
removal happens only at the end of the level) --
$\textit{parents}[\texttt{cog}]=[\texttt{dog},\texttt{log}]$.
Backtracking from \texttt{cog}: two branches, through \texttt{dog} and
through \texttt{log}, both leading back to \texttt{hit} -- exactly the
two shortest ladders the code prints.

\begin{complexitybox}
\textbf{Time Complexity.} $O(N \cdot 26L^2)$ for the level-by-level BFS
(same bound as Word Ladder I), plus time proportional to the total
length of all reconstructed paths for the backtracking step.
\textbf{Space Complexity.} $O(N \cdot L)$ for the parent map and word
sets, plus the space needed to store every returned path.
\end{complexitybox}

\begin{takeawaybox}
This chapter closes on the same lesson Print LIS and Print LCS taught
earlier in this book, now applied to an implicit graph: finding a
shortest distance and enumerating \emph{every} path that achieves it
are different problems, and the second one is solved by recording
parent relationships \emph{during} the same traversal that finds the
distance -- with one added wrinkle here, that ties at the same BFS
level require deferring when those parent relationships are finalized.
\end{takeawaybox}
